Index calculus is the method that should exploit the most structure in an
elliptic-curve discrete-logarithm instance, and in practice it has never
beaten Pollard rho on a prime-order curve of cryptographic shape. This
paper surveys where that question stands and measures it. We review the
literature by field type---extension fields (Gaudry, Diem, FGHR), prime-degree
binary fields and the contested first-fall-degree heuristics behind the
subexponential claims (Petit--Quisquater, Semaev, and the evidence of Huang,
Kosters and Yeo, Shantz--Teske and Galbraith--Gebregiyorgis), subfield curves
(Galbraith--Granger--Merz--Petit) and prime fields (Petit--Kosters--Messeng,
Amadori--Pintore--Sala, McGuire--Mueller)---and the generic lower bounds that
frame them.

We then measure the frontier with one native harness that charges rho,
baby-step giant-step, the kangaroo and several index-calculus pipelines in a
single counted unit on one-target workloads with sealed, replayable records.
Almost no cost in such records is complete, the rho baselines included, so we
price every arm's unpriced native work alike across a band of weights and rank
two methods only where the order holds across the band. On fixed ladders of
eight Koblitz curves ($\log_2 r$ from 11 to 39), six NIST-shaped prime curves
(20--40 bits) and five Kummer-tower curves---4{,}848 executions on public
targets, parameters chosen on separate training targets, every session
replay-audited---no index-calculus arm beats generic rho on any curve at any
weight. The closest, a compact Frobenius-orbit factor base, is $3.8$--$15\times$
signed-Frobenius rho with native work free; two-summand prime-field index
calculus goes from $10.5\times$ to $102\times$ rho between 20 and 40 bits; no
fitted cost exponent of a campaign index-calculus arm lies below rho's.

The measurements come from a fifteen-week engineering programme that we also
document: root-finding decomposition oracles (a known mechanism; its root
identity, Proposition~\ref{prop:s3}, is proved in Lean, and the oracle is
engineered to $383$--$1{,}535\times$ over matrix F4 per decomposition),
an $\FF_p$ F4 brought to msolve parity on square systems, FGHR's $D_4$
symmetry and its weighted and Frobenius-reduced forms cutting one $k=4$
decomposition by $48\times$ and then $370$--$481\times$, and two families
whose fitted cost grows more slowly than rho's: genus-three index calculus
after the cubic GHS transfer (about $q^{1.4}$ against $q^{3/2}$; $1.15\times$
a measured negation-map rho at $q=2^{20}$, down from $2.6\times$, with two
further rounds since) and $k=4$ index calculus over $\FF_{(2^m)^4}$, whose
gap to rho closes from $5{,}080\times$ to $6.1\times$ between $n\approx2^{24}$
and $2^{56}$ at Gaudry's $n^{-1/4}$, with a crossover extrapolated near
$2^{68}$. We also record the negative results, including a
conservation lemma for isogeny transfers and an equal-advice audit that
retires the programme's online wins over rho: Bernstein--Lange
preprocessing rho with the same table is $53$--$317\times$ faster, and at
$n=83$ it wins every pair matched on set-up time or memory. We close with open problems, each
stated as the experiment or proof that would settle it. Nothing here concerns a
deployed curve or key recovery.
